%! Unit 3: Quadratic Functions — Summative Assessment — Practice Test --- Study Copy (KEY)
\documentclass[10pt]{article}
\usepackage{algebra2-article}
\usepackage{algebra2-key}   % pulls in -boxes; do NOT also load -boxes
% Coordinate grid: {xmin}{xmax}{ymin}{ymax}
\newcommand{\lgrid}[4]{%
  \draw[step=1, linegray!40, very thin] (#1,#3) grid (#2,#4);
  \draw[->, semithick, charcoal] (#1,0)--(#2,0) node[below right, font=\scriptsize]{$x$};
  \draw[->, semithick, charcoal] (0,#3)--(0,#4) node[left, font=\scriptsize]{$y$};}
\newcommand{\dpt}[2]{\fill[charcoal] (#1,#2) circle (3pt);}
\newcommand{\vpt}[2]{\fill[navy] (#1,#2) circle (3.5pt);}

% Part-divider strip (headlinebox takes one arg: the background color)
\newcommand{\parthead}[1]{%
  \vspace{6pt}\noindent
  \begin{headlinebox}{forest}\color{white}\bfseries #1\end{headlinebox}%
  \vspace{4pt}}

\begin{document}

\pageheader{Unit 3: Quadratic Functions}{Practice Test --- Study Copy --- Answer Key}
\namedateperiod

\vspace{0.10in}
\begin{remindbox}
\small
\textbf{This is a practice test.} It mirrors the real Unit 3 test in length, format,
and the ideas it covers, but the numbers and contexts differ. Work every problem
with no notes, then check your answers against the key.
\end{remindbox}

% ======================================================================
\parthead{Part A --- Vocabulary (8 pts)}
Write the letter of the best description next to each term.

\vspace{4pt}
\noindent
\begin{minipage}[t]{0.42\textwidth}
\begin{enumerate}[leftmargin=2.2em, itemsep=7pt, label=\arabic*.]
  \item Vertex \ans{C}
  \item Axis of symmetry \ans{F}
  \item Zero (root) \ans{E}
  \item Discriminant \ans{B}
  \item Imaginary unit \ans{D}
  \item Completing the square \ans{H}
  \item Parabola \ans{A}
  \item Standard form \ans{G}
\end{enumerate}
\end{minipage}%
\hfill
\begin{minipage}[t]{0.55\textwidth}
\small
\begin{description}[leftmargin=1.4em, itemsep=3pt]
  \item[(A)] The U-shaped graph of a quadratic function.
  \item[(B)] The expression $b^2-4ac$; its sign tells the number and type of solutions.
  \item[(C)] The turning point of a parabola, where it reaches its maximum or minimum.
  \item[(D)] The number $i=\sqrt{-1}$, which satisfies $i^2=-1$.
  \item[(E)] An $x$-value where $f(x)=0$; an $x$-intercept of the parabola.
  \item[(F)] The vertical line $x=h$ that splits a parabola into mirror halves.
  \item[(G)] A quadratic written as $ax^2+bx+c$.
  \item[(H)] Rewriting $x^2+bx$ as a perfect-square trinomial by adding $\left(\tfrac{b}{2}\right)^2$.
\end{description}
\end{minipage}

% ======================================================================
\parthead{Part B --- Multiple Choice (12 pts)}
Circle the best answer.

\begin{enumerate}[leftmargin=1.9em, itemsep=9pt]
  \item The vertex of $y=(x-4)^2-1$ is
    \hfill \textbf{(A) $(4,-1)$} \ans{$\leftarrow$} \quad (B) $(-4,-1)$ \quad (C) $(4,1)$ \quad (D) $(-4,1)$
  \item $\sqrt{-25}=$
    \hfill \textbf{(A) $5i$} \ans{$\leftarrow$} \quad (B) $-5$ \quad (C) $25i$ \quad (D) $5$
  \item The discriminant of $x^2-4x+7=0$ is $-12$, so the equation has
    \hfill (A) two real solutions \quad (B) one real solution \quad \textbf{(C) two complex (non-real) solutions} \ans{$\leftarrow$} \quad (D) no solutions of any kind
  \item Factored completely, $x^2-49$ is
    \hfill \textbf{(A) $(x-7)(x+7)$} \ans{$\leftarrow$} \quad (B) $(x-7)^2$ \quad (C) $(x+7)^2$ \quad (D) prime
  \item The axis of symmetry of $y=x^2-6x+5$ is
    \hfill \textbf{(A) $x=3$} \ans{$\leftarrow$} \quad (B) $x=-3$ \quad (C) $x=6$ \quad (D) $x=5$
  \item A line and a parabola can intersect in \emph{at most} how many points?
    \hfill (A) $1$ \quad \textbf{(B) $2$} \ans{$\leftarrow$} \quad (C) $3$ \quad (D) infinitely many
\end{enumerate}

\begin{teachernote}
Item 2: pull the $i$ out first --- $\sqrt{-25}=\sqrt{25}\,\sqrt{-1}=5i$ (choice C forgets to
simplify $\sqrt{25}$). Item 3: $b^2-4ac<0\Rightarrow$ the $\pm\sqrt{\ }$ is imaginary, so two
complex conjugate roots and \emph{zero} $x$-intercepts (D is wrong --- there \emph{are} solutions,
just not real). Item 5: $x=-\tfrac{b}{2a}=\tfrac{6}{2}=3$.
\end{teachernote}

% ======================================================================
\parthead{Part C --- Short Answer \& Computation (40 pts)}
Show your work.

\begin{enumerate}[leftmargin=1.9em, itemsep=10pt]
  \item \textbf{Read the graph.} The parabola $f$ is shown below.
        \begin{center}
        \begin{tikzpicture}[scale=0.5]
          \lgrid{-3}{5}{-5}{5}
          \draw[very thick, forest, smooth, samples=100, domain=-2:4] plot (\x, {(\x-1)*(\x-1)-4});
          \vpt{1}{-4}
          \dpt{-1}{0}
          \dpt{3}{0}
          \dpt{0}{-3}
        \end{tikzpicture}
        \end{center}
        Vertex: \ans{$(1,-4)$}; \quad axis of symmetry: \ans{$x=1$}; \quad $y$-intercept: \ans{$(0,-3)$};
        \\[4pt]
        zeros: \ans{$x=-1,\ x=3$}; \quad minimum \emph{or} maximum value: \ans{minimum $-4$}; \quad range: \ans{$y\ge -4$}

  \item \textbf{From standard form.} For $f(x)=x^2-6x+5$, use $x=-\dfrac{b}{2a}$ to find the vertex,
        give the axis of symmetry, and state whether the parabola opens up or down.
        \ans{$x=-\dfrac{-6}{2(1)}=3$; $f(3)=9-18+5=-4$, so the vertex is $(3,-4)$. Axis of symmetry $x=3$. Since $a=1>0$, it opens \emph{up} (so $-4$ is a minimum).}

  \item \textbf{Solve by factoring.}
        \[ x^2+2x-15=0 \]
        \ans{$(x+5)(x-3)=0\Rightarrow x=-5$ or $x=3$. Check: $(-5)^2+2(-5)-15=0$ \checkmark, $3^2+2(3)-15=0$ \checkmark.}

  \item \textbf{Solve using square roots.}
        \[ (x-2)^2=49 \]
        \ans{$x-2=\pm 7\Rightarrow x=2+7=9$ or $x=2-7=-5$.}

  \item \textbf{Solve by completing the square.}
        \[ x^2+6x-7=0 \]
        \ans{$x^2+6x=7$; add $\left(\tfrac{6}{2}\right)^2=9$ to both sides: $(x+3)^2=16$; $x+3=\pm 4$, so $x=1$ or $x=-7$.}

  \item \textbf{Complex numbers.} Let $z_1=3+2i$ and $z_2=1-4i$. Write each answer in $a+bi$ form.
        \\[4pt]
        (a) $z_1+z_2=\ans{$4-2i$}$ \qquad (b) $z_1-z_2=\ans{$2+6i$}$
        \\[8pt]
        (c) $z_1\cdot z_2=\ans{$(3+2i)(1-4i)=3-12i+2i-8i^2=3-10i+8=11-10i$}$

  \item \textbf{Quadratic formula \& discriminant.} Solve $2x^2+3x-2=0$. First compute the
        discriminant and state how many real solutions to expect, then solve.
        \ans{$b^2-4ac=3^2-4(2)(-2)=9+16=25>0$, so \emph{two real} solutions. $x=\dfrac{-3\pm\sqrt{25}}{2(2)}=\dfrac{-3\pm 5}{4}$, giving $x=\tfrac12$ or $x=-2$.}

  \item \textbf{System.} Solve the linear--quadratic system algebraically. Give every solution as an
        ordered pair.
        \[ y=x^2-2x-3 \qquad y=x-3 \]
        \ans{Set equal: $x^2-2x-3=x-3\Rightarrow x^2-3x=0\Rightarrow x(x-3)=0$, so $x=0$ or $x=3$. Then $y=x-3$: $x=0\Rightarrow(0,-3)$; $x=3\Rightarrow(3,0)$. Two solutions: $(0,-3)$ and $(3,0)$.}
\end{enumerate}

\begin{teachernote}
Item 1: accept ``opens up / minimum'' language; range is $y\ge-4$ (closed, vertex included).
Item 6c: full credit requires distributing all four terms \emph{and} replacing $i^2=-1$ ---
watch for students who leave $-8i^2$. Item 7: since the discriminant is a perfect square ($25$),
the roots are rational; a non-perfect-square discriminant would leave a radical.
\end{teachernote}

% ======================================================================
\parthead{Part D --- Extended Response (12 pts)}
Explain your reasoning in full sentences.

\begin{enumerate}[leftmargin=1.9em, itemsep=16pt]
  \item \textbf{Projectile.} A ball is thrown upward from a rooftop. Its height (in feet) after
        $t$ seconds is
        \[ h(t)=-16t^2+32t+48. \]
        \textbf{(a)} From which feature of the model do you read the ball's \emph{starting height}?
        Give that height.
        \textbf{(b)} Find the \emph{maximum} height and the time it occurs. Which feature of the graph
        is this?
        \textbf{(c)} When does the ball hit the ground? Which feature answers this, and why do you
        reject the other solution?
        \ans{(a) The starting height is the $y$-intercept, $h(0)=48$ feet. (b) The maximum is the \emph{vertex}. $t=-\dfrac{b}{2a}=-\dfrac{32}{2(-16)}=1$ second; $h(1)=-16+32+48=64$ feet. So the ball reaches a maximum height of $64$ ft at $t=1$ s. (c) The ground is $h=0$ --- the \emph{zeros}. $-16t^2+32t+48=0\Rightarrow t^2-2t-3=0\Rightarrow(t-3)(t+1)=0$, so $t=3$ or $t=-1$. We reject $t=-1$ because a negative time is meaningless here; the ball lands at $t=3$ seconds.}

  \item \textbf{Discriminant reasoning.} Consider the equation $x^2-4x+k=0$, where $k$ is a constant.
        \textbf{(a)} Write the discriminant in terms of $k$.
        \textbf{(b)} For what values of $k$ does the equation have two real solutions? one? two complex
        solutions?
        \textbf{(c)} Explain how the discriminant is connected to the number of $x$-intercepts of the
        parabola $y=x^2-4x+k$.
        \ans{(a) $b^2-4ac=(-4)^2-4(1)(k)=16-4k$. (b) Two real: $16-4k>0\Rightarrow k<4$. One (double) real: $16-4k=0\Rightarrow k=4$. Two complex: $16-4k<0\Rightarrow k>4$. (c) The real solutions are exactly the $x$-values where the parabola meets the $x$-axis. So a positive discriminant ($k<4$) means two $x$-intercepts, a zero discriminant ($k=4$) means one (the vertex sits on the $x$-axis), and a negative discriminant ($k>4$) means no $x$-intercepts (the whole parabola stays above the axis).}
\end{enumerate}

\begin{teachernote}
\textbf{Scoring Part D (6 pts each).} D1: 2 pts (a) $y$-intercept named \emph{and} $48$ ft;
2 pts (b) vertex $t=1$, height $64$ ft; 2 pts (c) zeros, $t=3$, with the $t=-1$ rejection reasoned
from context. D2: 1 pt (a) $16-4k$; 3 pts (b) all three cases with correct boundary $k=4$; 2 pts
(c) ties discriminant sign to the \emph{count of $x$-intercepts}. Do not give full (c) credit for
restating (b) without the graph connection.
\end{teachernote}

\end{document}
